\begin{abstract}
    Alignment research addresses three structural layers: specification
    (what to optimize), generalization (whether the agent reliably
    pursues it), and verification (whether pursuit is measured
    faithfully).  Most research targets the first two layers.  We argue
    that verification is the prerequisite layer: when verification fails,
    solutions to specification and generalization do not bind.  The
    verification problem has a precise form: a \emph{write-access
    asymmetry} in which the agent's action space includes actions that
    modify the evaluator's inputs.  This asymmetry grows
    monotonically with capability.  We formalize \emph{substrate-exclusive
    witnesses}: verification processes rooted in substrates the claimant
    cannot write to, providing the only structural barrier to evaluator
    gaming that is invariant under capability scaling.  A cryptographic
    commitment protocol (binding, hiding; with an optional
    zero-knowledge variant for sensitive claims) enforces the
    substrate-exclusivity barrier without requiring trust, and an
    append-only verification ledger provides the persistent record.
    Human judgment enters at exactly one layer (benchmark admissibility
    via a governance fork mechanism) while per-claim verification is
    fully algorithmic at machine speed.  Within the goal-frontier
    maximization (GFM) framework, the construction closes the
    observation-channel integrity assumption underlying multi-channel
    attribution, upgrades SPRT-based scorpion detection from estimated
    to committed baselines, and yields the \emph{anti-monopolar
    verification theorem}: a monopolar population ($m = 1$) loses not
    only the safety guarantee but the ability to detect that it has
    lost it.
\end{abstract}
